% Generated draft: requires mathematical and source review.
\documentclass[11pt]{article}
\usepackage[margin=25mm]{geometry}
\usepackage{fontspec}
\setmainfont{Latin Modern Roman}
\usepackage{amsmath,amssymb,mathtools,unicode-math}
\setmathfont{Latin Modern Math}
\usepackage{enumitem,needspace,xurl,hyperref}
\setlength{\parindent}{0pt}
\setlength{\parskip}{5pt}
\setlength{\emergencystretch}{3em}
\newcommand{\N}{\mathbb{N}}
\newcommand{\Z}{\mathbb{Z}}
\newcommand{\Q}{\mathbb{Q}}
\newcommand{\R}{\mathbb{R}}
\newcommand{\C}{\mathbb{C}}
\newcommand{\F}{\mathbb{F}}
\newcommand{\A}{\mathbb{A}}
\newcommand{\PP}{\mathbb{P}}
\newcommand{\E}{\mathbb{E}}
\newcommand{\eps}{\varepsilon}
\newcommand{\dd}{\,\mathrm{d}}
\begin{document}
\section*{UnsolvedMath ID: 11000284}
\textbf{Source problem code:} catalogue entry
\subsection{Definitions and mathematical statement}
Let \(F_n\) be a free group of rank \(n\), and let \(G\leq\operatorname{Out}(F_n)\). A group is virtually abelian if it contains an abelian subgroup of finite index. In the ordinary bounded-cohomology convention, \(H_b^2(G;\R)\) is the cohomology of bounded real group cochains with trivial coefficients, without silently replacing it by ordinary cohomology or a reduced quotient.

\emph{Question 3.4.} If \(G\) is not virtually abelian, must
\[
 \dim_\R H_b^2(G;\R)=\infty?
\]
The source imposes no finite-generation, finite-index, or special geometric hypothesis on \(G\), and specifies degree \(2\) and real coefficients. The allowed free-group rank range is implicit.

The following context discusses embeddings as rank increases. It notes the evident embeddings \(\operatorname{GL}_n(\Z)\to\operatorname{GL}_m(\Z)\) and \(\operatorname{Aut}(F_n)\to\operatorname{Aut}(F_m)\) for \(m\geq n\), while saying no comparable embedding of outer automorphism groups is obvious. It credits Bogopolski and Puga [10] with an embedding for a special rank relation, printed in the extraction as
\[
 m=1+(n-1)kn,\qquad \gcd(k,n-1)=1,
\]
with \(k\) a natural number. The exponent placement is unresolved: the compressed \(kn\) may mean the product \(kn\) or \(k^n\). The source also describes restriction to a suitable characteristic subgroup ``of \(F_m\),'' but does not explain how that restriction realizes the displayed direction. Both rank and characteristic-subgroup conventions need source verification; no corrected embedding theorem is silently asserted. This adjacent embedding discussion supplies no proof of the bounded-cohomology question.
\subsection{Short English statement}
Must every non-virtually-abelian subgroup of \(\operatorname{Out}(F_n)\) have infinite-dimensional real bounded cohomology in degree two? The nearby rank-embedding claim is retained with unresolved exponent and subgroup-index notation.
\subsection{Sources}
\begin{itemize}
\item Local UnsolvedMath registry, permanent ID 11000284, source code catalogue entry, Question 3.4, PDF page 333. No individual proposer is supplied.
\item Farb, \emph{Problems on Mapping Class Groups and Related Topics}. Supplied URL: \url{https://www.math.uchicago.edu/~farb/papers/mcgbook.pdf}; recorded access 2026-07-29; extraction pdf-font-structure. The volume contains 329 explicitly typeset headings across chapters, with repeated numbering disambiguated by page.
\item Bogopolski and Puga [10] are cited for the special-rank embedding, but the full entry is absent. The summary mentions Bestvina--Fujiwara-type results without a precise theorem or complete reference.
\item Supplied status open and OPEN-TRIAGE, with summary language ``partial'' and ``Cannot verify.'' Current primary-literature status and redistribution terms for the publicly accessible author/publisher copy remain NEEDS\_REVIEW.
\end{itemize}
\end{document}
